\section{A finite-chain clock comparison}
\label{app:clock-comparison}

This appendix proves the bounded-array clock lemma stated in
Section~\ref{sec:model}. Both cutoff arguments use it to pass from
shifted Poisson averages to a binomial average after extracting a fixed
holding probability from the update kernel. The parameters and arrays
below are those of Lemma~\ref{lem:input-clock}; the proof uses no
property of the SK model. It follows the shifted-clock comparison of
Chen and Saloff-Coste~\cite[Theorem~1.1 and Lemmas~3.1--3.2]{ChenSaloffCoste2013},
with the finite approximation fixed before the arrays are tested.
The Gaussian approximation below uses signed coefficients: cancellation
permits translates of the wider Poisson Gaussian to approximate the
narrower binomial Gaussian.

\begin{proof}[Proof of Lemma~\ref{lem:input-clock}]
The proof has three steps.  We approximate the two rescaled clocks in
\(L^1\) by Gaussian densities, approximate the narrower Gaussian by a
finite signed sum of translates of the wider one, and then test this
single finite approximation against the bounded arrays.  Fixing that
approximation before taking the limit will also give the uniform statement.

\emph{Step 1: the two local limits.}
For an integer
$M$ and $j\in\mathbb Z$, let
\[
 I_{M,j}=\left[\frac{j-q M}{\sqrt M},
                   \frac{j+1-q M}{\sqrt M}\right).
\]
The histogram density of an integer-valued random variable $Z$ is
$\sqrt M\,\P(Z=j)$ on $I_{M,j}$.  Denote by $p_{M,y}$ the histogram
for $\operatorname{Poisson}(q M+y\sqrt M)$ and by $q_M$ the
histogram for $\operatorname{Binomial}(M,q)$, assigning zero
probability outside their supports.  For $v>0$ write
\[
 \varphi_v(x)=\frac{1}{\sqrt{2\pi v}}e^{-x^2/(2v)}.
\]
We claim that for every fixed $y$,
\begin{equation}
 \|p_{M,y}-\varphi_q(\,\cdot-y)\|_1\longrightarrow0,
 \qquad
 \|q_M-\varphi_{q(1-q)}\|_1\longrightarrow0.
 \label{eq:inputclock-local-limits}
\end{equation}

Here is a proof from probability ratios.  Let $\pi_M(j)$ denote the
probability mass function of either variable, let $k_M$ be an integer
nearest its mean, and put $v=q$ in the Poisson case and
$v=q(1-q)$ in the binomial case.  Their successive ratios
are, respectively,
\[
 \frac{m}{j+1},\qquad
 \frac{(M-j)q}{(j+1)(1-q)},
 \qquad m=q M+y\sqrt M.
\]
For every fixed $H<\infty$, expanding their logarithms at $k_M$
gives, uniformly for integers $|\ell|\le H\sqrt M$,
\[
 \log\frac{\pi_M(k_M+\ell+1)}{\pi_M(k_M+\ell)}
 =-\frac{\ell}{vM}+O_H(M^{-1}).
\]
The constants here may depend on the fixed $q$ and $y$.  Indeed,
the deviations of the numerator and denominator from their central
values are $O_H(\sqrt M)$, and the remainder in
$\log(1+h)=h+O(h^2)$ is $O_H(M^{-1})$; rounding the mean costs the
same order.  All these indices belong to the supports for large $M$.
Summing forward or backward from $k_M$ proves
\begin{equation}
 \frac{\pi_M(k_M+\ell)}{\pi_M(k_M)}
 =\exp\left(-\frac{\ell^2}{2vM}+O_H(M^{-1/2})\right)
 \qquad(|\ell|\le H\sqrt M).
 \label{eq:inputclock-ratio-expansion}
\end{equation}

The ratio estimate determines the shape near the mean.  Its
normalization remains to be identified.  Set $A_M=\sqrt M\,\pi_M(k_M)$ and
$J_H=\int_{-H}^H e^{-x^2/(2v)}\,dx$.
Equation~\eqref{eq:inputclock-ratio-expansion} and Riemann sums imply
\[
 \P(|Z-k_M|\le H\sqrt M)=A_M(J_H+o(1)).
\]
Taking, for example, $H=1$ shows that $A_M$ is bounded.
The variance of $Z$ divided by $M$ tends to $v$, and its mean is within
$1/2$ of $k_M$.  Chebyshev's inequality therefore gives
\[
 \limsup_{M\to\infty}\P(|Z-k_M|>H\sqrt M)\le\frac{v}{H^2}.
\]
For fixed $H$ this bounds the limit inferior and limit superior of
$A_M$ by $(1-v/H^2)/J_H$ and $1/J_H$, respectively.  Letting
$H\to\infty$ yields
$A_M\to(2\pi v)^{-1/2}$.
Equation~\eqref{eq:inputclock-ratio-expansion} now proves uniform convergence
of the histogram densities on every compact interval, with limiting
center $y$ in the Poisson case and zero in the binomial case.
The same Chebyshev estimate controls their tails; the histogram
intervals move an observation by at most $M^{-1/2}$ in the rescaled
coordinate.  Gaussian tails then give the whole-line $L^1(\R)$ convergence
in Equation~\eqref{eq:inputclock-local-limits}.

\emph{Step 2: approximation by translates.}
We claim that the narrower Gaussian belongs to the closed linear span,
in $L^1$, of translates of the wider Gaussian.  Every derivative
$\varphi_q^{(k)}$ belongs to that closed span: its finite
differences are finite linear combinations of translates, and their
quotients converge to the derivative in $L^1$.  For completeness,
iterating the Fundamental Theorem of Calculus expresses the
$k$th forward-difference quotient as
\[
 \int_{[0,1]^k}\varphi_q^{(k)}
           (x+h(t_1+\cdots+t_k))\,dt_1\cdots dt_k.
\]
Gaussian derivatives are integrable polynomial multiples of a
Gaussian, and their translations are continuous in $L^1$ by dominated
convergence.  This proves the asserted limit.  Repeated differentiation
also gives $\varphi_q^{(k)}=p_k\varphi_q$, where $p_k$ is
a polynomial of degree $k$ with leading coefficient $(-1/q)^k$.
The nonzero leading coefficients therefore let us express every
polynomial multiple of $\varphi_q$ as a finite linear combination of
these derivatives.  Such multiples lie in the same closed span.

Put $v=q$ and $w=q(1-q)$.  The ratio of the two
densities is
\[
 \frac{\varphi_w(x)}{\varphi_v(x)}
 =\sqrt{\frac vw}\,e^{-c x^2},
 \qquad c=\frac12\left(\frac1w-\frac1v\right)>0.
\]
If $V$ is centered Gaussian of variance $2c$, then
$e^{-cx^2}=\E\cos(Vx)$.  This identity follows directly from Gaussian
integration by parts: the characteristic function has value one at
zero and derivative $-2cx$ times itself.  Truncate the expectation to
$|V|\le R$.  The omitted contribution has absolute value at most
$\P(|V|>R)$, uniformly in $x$.  For fixed $R$, replace cosine by its
Taylor polynomials.  On $|V|\le R$, their absolute values are bounded
by $e^{R|x|}$, which is integrable against $\varphi_v(x)\,dx$.
Dominated convergence thus approximates the truncated expectation by
polynomials in $L^1(\varphi_v(x)\,dx)$.  Sending $R$ to infinity and
using the preceding derivative argument proves the claimed closed-span
property for $\varphi_w$.

\emph{Step 3: test one finite approximation against the arrays.}
Extend $a_{n,j}$ by zero to $j<0$, and let $U_n$ be the
function equal to $a_{n,j}$ on $I_{M_n,j}$.  If $B=0$ the result is
immediate, so assume $B>0$.  Given $\epsilon_*>0$, choose one finite
signed sum
\[
 F=\sum_{\ell=1}^L a_\ell\varphi_q(\,\cdot-y_\ell),
 \qquad
 \|\varphi_{q(1-q)}-F\|_1
                         <\frac{\epsilon_*}{4B}.
\]
The number of terms, the real shifts, and the real coefficients are
fixed before $n$ tends to infinity.  Set
$A_0=\max\{1,\sum_\ell|a_\ell|\}$ and abbreviate
$\psi=\varphi_{q(1-q)}$,
$\phi_\ell=\varphi_q(\,\cdot-y_\ell)$.
Then
\begin{align}
 \left|\int q_{M_n}U_n\right|
 &\le B\|q_{M_n}-\psi\|_1+B\|\psi-F\|_1 \notag\\
 &\quad+\sum_{\ell=1}^L|a_\ell|
       \left(\left|\int p_{M_n,y_\ell}U_n\right|
             +B\|p_{M_n,y_\ell}-\phi_\ell\|_1\right).
 \label{eq:inputclock-finite-approximation}
\end{align}
For large $n$, the first $L^1$ error is below $\epsilon_*/(4B)$,
each of the finitely many other $L^1$ errors is below
$\epsilon_*/(4BA_0)$, and each of the Poisson averages is below
$\epsilon_*/(4A_0)$.  The right side is then at most $\epsilon_*$.
This proves the result.  The same bound proves its uniform version:
the finite approximation is common to the family, the local-limit
errors depend only on $M$ and the fixed shifts, and the hypothesis
controls the finite list of Poisson averages uniformly.
\end{proof}
